\section{Evaluation}

In this chapter we evaluate the rstrategies library in two general directions: runtime performance and applicability to different programming languages.
We evaluate the runtime performance by comparing the results of benchmarks executed with and without storage strategies.
We show the generality of the library by demonstrating two programming languages of very different nature that have been successfully adapted to using rstrategies.







\subsection{Performance}
\todo{Does the performance benefit reported in the PyPy paper carry over?}

In this section we present the results of experiments measuring the execution time of benchmarks to evaluate the performance benefits of storage strategies.
The experiments have been conducted on a 64-bit Windows 8.1 machine with 8GB of RAM and an Intel Core i7-2620M CPU with two physical and four virtual cores, clocked at up to 3.40GHz during single-core execution.
During the experiments, no other user applications were running and the experiment process was given increased priority in order to minimize influences from operating system processes.
We used the RSqueak VM at commit 1060:7a0ce39abc12 (\date{16.10.2014}), Pypy at commit 72481:35fdf446e439 (\date{24.07.2014}) and a Squeak 4.5 image to take the measurements.

To evaluate the runtime performance of rstrategies, all experiments were conducted twice.
First, all specialized storage strategies have been disabled via a command line switch.
In this setup, the \lstinline{ListStrategy} is used by every object or collection created during the benchmarks.
In the second setup, specialized storage strategies were enabled.
In the case of the RSqueak VM, these strategies are \lstinline{AllNilStrategy}, \lstinline{SmallIntegerOrNilStrategy} and \lstinline{FloatOrNilStrategy}.
We compare the runtime performance of these two experimental setups to evaluate the impact of the specialized storage strategies.

The following list gives short explanations of the seven benchmarks executed during the experiments.
These are all common benchmarks based on The Computer Language Benchmarks Game \footnote{\url{http://benchmarksgame.alioth.debian.org/}}.
\begin{description}
\item[AStar]{The AStar path searching algorithm is used to find a solution through two predefined mazes. This benchmark creates a big graph of interconnected objects and generates lots of non-recursive messages. Two different mazes are solved, called AStar1 and AStar2; the second maze is larger.}
\item[BinaryTree]{Balanced binary trees of a given depth are constructed and then walked, summing up the elements of all leaf-nodes. This benchmark generates many recursive messages.}
\item[Blowfish]{The symmetric-key block cipher Blowfish is used to encrypt and decrypt fixed challenge data. This benchmark is heavy on basic number arithmetic, very few objects are created.}
\item[DeltaBlue]{DeltaBlue is a constraint solver benchmark. Different constraints are created for a number of variables, then the solver finds optimal solutions for the given constraints. This benchmark is heavy on message sends and conditional logic.}
\item[NBody]{The NBody algorithm creates a number of objects with an assigned movement vector and mass, then performs a round-based simulation of their movement and interactions using Newton's laws of motion. This benchmark balances basic arithmetics and message sends.}
\item[Richards]{The Richards benchmarks simulates the task scheduler of an operating system kernel. Multiple classes of tasks can be created with different priorities, then the tasks are scheduled until every task has finished. The main workload of this benchmark consists of message sends and conditional logic.}
\item[SplayTree]{This benchmark creates a splay tree of a fixed size and then performs a number of tree modifications by removing the greatest key node in the tree. This benchmark creates and modifies a big interconnected graph of objects.}
\end{description}

For all experiments, entire benchmark suite is executed within the same VM instance.
This creates a more realistic setup, since a real-world application is likely to be long-running and to perform multiple different tasks.
Appendix \ref{app:benchmarksSeparateVMs} contains the benchmark results when every benchmark is executed in its own VM instance.

\subsubsection{Results}

The execution time of the benchmarks was measured by recording timestamps at the beginning and end of every benchmark.
Every benchmark was executed 50 times to achieve a satisfactory significance level.
From the 50 measurements taken for every benchmark, the first \emph{n} are excluded from the analysis to eliminate the impact of the JIT warmup time.
The number \emph{n} is determined manually for every benchmark, since the time it takes to reach steady performance differs from one benchmark to another.
Table \ref{tab:excluded} lists the number of excluded measurements for every benchmark.

\begin{table}
\centering
\begin{tabular}[ht]{ l c@{\ \ \ \ }c }
\toprule
Benchmark & \parbox{2.4cm}{Excluded\\Measurements} & \parbox{2.4cm}{Analysed\\Measurements} \\
\midrule
AStar1 & 14 & 36 \\
AStar2 & 14 & 36 \\
BinaryTree & 3 & 47 \\
BlowfishDecryption & 13 & 37 \\
BlowfishEncryption & 16 & 34 \\
DeltaBlue & 8 & 42 \\
NBody & 3 & 47 \\
Richards & 10 & 40 \\
SplayTree & 12 & 38 \\
\bottomrule
\end{tabular}
\caption{Excluded and analysed measurements for the different benchmarks.}
\label{tab:excluded}
\end{table}

The remaining measurements have been summarized by means of confidence intervals using a 5\% significance level.
The confidence intervals comparing the two VMs have been computed following the procedure proposed by \citeauthor{rigorousBenchmarking} \cite{rigorousBenchmarking}.
The confidence intervals summarizing the results of the individual VMs show the standard error of the population, assuming normal distribution.
Table \ref{tab:runtimeResults} summarizes the results, each line represents one of the benchmarks.
The table is divided in three columns, each summarizing different results for the respective benchmark.
The first column shows the results for the RSqueak VM \emph{without} storage strategies, i.e. with the rstrategies library disabled.
The second column shows the results for the RSqueak VM with storage strategies enabled.
The numbers in the first two columns are given in milliseconds.
The third column shows the relative performance change from the first to the second section.
A positive number indicates increased performance, meaning that the benchmark was executed in less time.

\begin{table}
\centering
\begin{tabular}[ht]{%
@{}l%
r@{\,\si{\milli\second}}>{\ \smaller\ensuremath{\pm}}r%
r@{\,\si{\milli\second}}>{\ \smaller\ensuremath{\pm}}r%
r@{\,\si{\percent}}>{\ \smaller\ensuremath{\pm}}r@{}}
\toprule
Benchmark & \multicolumn{2}{l}{\parbox{2cm}{Without\\Strategies}} & \multicolumn{2}{l}{\parbox{2cm}{With\\Strategies}} & \multicolumn{2}{l}{\parbox{2cm}{Performance\\Change}} \\
\midrule
AStar1 & 80.64 & 1.92 & 55.81 & 2.61 & \(+\)46.93 & 6.58\\
AStar2 & 351.17 & 8.15 & 288.81 & 2.68 & \(+\)21.21 & 2.25\\
BinaryTree & 187.47 & 1.00 & 174.94 & 1.77 & \(+\)7.26 & 1.01\\
BlowfishDecryption & 422.16 & 2.05 & 429.16 & 2.49 & \(-\)1.62 & 0.64\\
BlowfishEncryption & 423.85 & 2.00 & 427.15 & 2.47 & \(-\)0.76 & 0.63\\
DeltaBlue & 99.86 & 2.24 & 98.31 & 2.16 & \(+\)1.57 & 2.62\\
NBody & 271.38 & 2.15 & 274.09 & 2.16 & \(-\)0.98 & 0.94\\
Richards & 165.97 & 2.80 & 162.95 & 4.14 & \(+\)2.11 & 2.29\\
SplayTree & 781.16 & 2.25 & 469.92 & 2.89 & \(+\)66.28 & 0.97\\
\bottomrule
\end{tabular}
\caption{Performance measurements for different benchmarks.
From left to right: RSqueak VM without specialized storage strategies, RSqueak VM with specialized storage strategies, relative performance change between the first two.
Measurements given as \SI{95}{\percent} confidence intervals.}
\label{tab:runtimeResults}
\end{table}

The AStar1 and SplayTree benchmarks showed considerable increase in performance, while Blowfish and NBody showed a minor decrease.
The other three benchmarks showed minor increases between 1\% and 7\%.
Since the means and standard errors shown in the first two columns differ quite substantially, it would not be meaningful to construct a single confidence interval from all benchmarks.
Instead, we observe that storage strategies provide high performance benefits for certain problems, while the potential performance penalty remains very low.
The average half-width of the confidence intervals is 1.99\%, which is reasonably low for a significance level of 5\%.

The two benchmarks that benefit most, AStar and SplayTree, both create large graphs of objects.
Therefore, it is not surprising that they benefit a lot from an allocation removal optimization.
The benchmarks with slight performance losses, Blowfish and NBody, are the two benchmarks with the highest focus on arithmetical operations.
Since there are not many object allocations happening here, we can measure the overhead of storage strategies as a slight performance drop.






\subsection{Memory}
\todo{Does the memory benefit carry over?}
%\begin{figure}
%\label{fig:memory}
%\caption{Memory overhead...}
%\missingfigure{\Huge Need the figure here!}
%\end{figure}





\subsection{Generality}
\todo{Can the framework be applied to other languages? Topaz, Pycket}






\subsection{Transitions}
\todo{Are the strategies we identified as most important really that?}
%\begin{figure}
%\label{fig:transitions}
%\caption{Transitions...}
%\missingfigure{\Huge Need the figure here!}
%\end{figure}
